% Propuesta separada; no modifica el bloque común del artículo I.
% Procedencia: artículo II, revisión 05, registro_incidencias.tex y
% registro_imagen_integral.tex; aclaraciones de índices de revisiones 06/07.
% Inserción prevista inmediatamente antes de subsec:k-terminal.
\subsection{Evaluación de incidencias y determinación de las coordenadas transversales}
\label{iii:r25:incidencias}

El registro dodecafásico conserva una historia de visitas y trazas
orientadas. Sus coordenadas transversales se obtienen mediante una
composición de operaciones: evaluación aritmética de las visitas,
recuento por ventanas, normalización decimal y diferencias cíclicas.
La reconstrucción del registro a partir de esas diferencias constituye
otra operación, cuya existencia y unicidad se demostrarán después.
Esta separación permite precisar tanto los datos utilizados por el
extractor como la información que transmite a la lectura armónica.

Sea $\mathscr L$ un registro de incidencias producido por la selección,
el transporte y la actualización de una historia APP--TRIT--TPK.
Cada visita conserva su ruta, posición APP, hoja, instante,
multiplicidad y clase de frontera; cada traza conserva sus miembros,
orientación, multiplicidad y longitud en la unidad declarada por su
lector. Las ventanas son
\[
 E_m=\{9(m-1)+1,\ldots,9m\},\qquad 1\leq m\leq12.
\]
La evaluación de la hoja de suma o producto determina el entero
$n(v)$ de cada visita y su división positiva
\[
 n(v)=r(v)+9q(v),\qquad 1\leq r(v)\leq9.
\]
El residuo se calcula sobre esa evaluación. La clase de frontera de
una visita de residuo $9$ se conserva mediante
$\epsilon_\partial(v)=1$ en la corona y $-1$ en el interior.

La presentación incidencial del registro utiliza los tres recuentos
\begin{align}
 A_m&=\sum_{v\in E_m}w(v)
       \mathbf1_{\{1,3,5,7\}}(r(v)),\label{iii:r25:recuento-a}\\
 C_m&=\sum_{j\geq0}
       \bigl(\nu^-_{120,m}(j)-\nu^+_{120,m}(j)\bigr),
       \label{iii:r25:recuento-c}\\
 V_m&=\sum_{v\in E_m}w(v)
       \mathbf1_{\{9\}}(r(v))\epsilon_\partial(v).
       \label{iii:r25:recuento-v}
\end{align}
Aquí $w(v)$ es la multiplicidad incidencial y
$\nu^\pm_{120,m}(j)$ cuenta las trazas orientadas del canal indicado
de longitud reducida $120(1+3j)$, según la enumeración del registro.
La familia de trazas y su multiplicidad son datos de la evaluación,
con la regla de producción que les corresponde. La suma entera requiere
soporte finito o una construcción que determine su significado.
La longitud reducida de una traza y el número de instantes de la
ventana son magnitudes tipadas diferentes; la unidad del lector
especifica su relación. Los símbolos $A_m,C_m,V_m$ designan aquí
recuentos enteros, distintos del par angular y del operador constitutivo.

Sea $[x]_{10}\in\{0,\ldots,9\}$ el representante euclídeo.
La normalización conserva también el cociente en
\[
 x=[x]_{10}+10\lfloor x/10\rfloor.
\]
Se forma el bloque
\begin{equation}
 z_m=100[A_m]_{10}+10[C_m]_{10}+[V_m]_{10},
 \qquad 0\leq z_m\leq999.
 \label{iii:r25:bloque}
\end{equation}
Con índices cíclicos módulo $12$, sean
\[
 (D_a z)_m=z_m-z_{m+a},\qquad a\in\{3,4\},
 \qquad Q(z)=\sum_{m=1}^{12}z_m.
\]
La composición incidencial queda expresada por
\begin{equation}
 \mathcal E_{\rm cnt}(\mathscr L)
 =\mathcal D(z):=(D_3z,D_4z,Q(z))\in\mathbb Z^{25}.
 \label{iii:r25:composicion}
\end{equation}
Sus argumentos son las visitas y trazas con sus reglas de recuento.
La evaluación no recibe $K$, $U$, $\alpha$ ni las coordenadas
transversales esperadas. Para identificarla con
$\mathcal E_{108}^{90,120}$ se conserva además la composición con
la familia y las incidencias efectivas del estado terminal. La
invertibilidad de $\mathcal D$ determina una reconstrucción posterior;
por sí sola no selecciona esa familia de incidencias.

\begin{proposition}[Información residual determinada por el extractor]
\label{iii:r25:residuos}
Sean $N,N'\in\mathbb Z^{12\times3}$ dos registros de recuentos,
ordenados por $(A,C,V)$. Para las operaciones
\eqref{iii:r25:bloque}--\eqref{iii:r25:composicion},
\[
 \mathcal E_{\rm cnt}(N)=\mathcal E_{\rm cnt}(N')
 \quad\Longleftrightarrow\quad N-N'\in10\mathbb Z^{36}.
\]
Los $36$ residuos sectoriales determinan exactamente la salida de
$25$ coordenadas. Los cocientes y las procedencias constituyen
información adicional del registro.
\end{proposition}
\begin{proof}
La igualdad módulo $10$ produce los mismos bloques y sus mismas
coordenadas transversales. Recíprocamente, la igualdad de las
coordenadas da $D_3(z-z')=D_4(z-z')=0$. Los pasos $3$ y $4$
generan el ciclo de longitud $12$, por lo que $z-z'$ es constante.
La igualdad de $Q$ anula esa constante. La expansión de cada entero
de $0$ a $999$ en tres cifras decimales es única; se recuperan así
los tres residuos de cada ventana.
\end{proof}

La proposición es válida para todo el dominio entero indicado. No
postula que las $36$ clases sean independientes sobre el subconjunto
de registros admisibles producido por TPK.

\subsection{Imagen integral e inversión de las coordenadas transversales}
\label{iii:r25:imagen-seccion}

Para expresar la inversión se reindexan los bloques por
$i\in\mathbb Z/12\mathbb Z$, de $0$ a $11$. Sean
$(Sz)_i=z_{i+1}$, $D_3=I-S^3$ y $D_4=I-S^4$.
Dada una terna entera $(b,c,q)$, defínanse
\begin{equation}
 w_i=c_i-b_{i+1},\qquad
 P_0=0,\quad P_i=\sum_{j=0}^{i-1}w_j\ (1\leq i\leq12),
 \qquad T(w)=\sum_{i=0}^{11}P_i.
 \label{iii:r25:incrementos}
\end{equation}

\begin{theorem}[Caracterización de la imagen integral]
\label{iii:r25:imagen}
Una terna $(b,c,q)\in\mathbb Z^{25}$ pertenece a
$\operatorname{im}\mathcal D$ si y sólo si
\begin{align}
 b_i&=w_i+w_{i+1}+w_{i+2},\qquad 0\leq i<12,
       \label{iii:r25:relaciones}\\
 \sum_{i=0}^{11}w_i&=0,\label{iii:r25:cierre}\\
 q+T(w)&\equiv0\pmod {12}.\label{iii:r25:congruencia}
\end{align}
Su única preimagen entera es
\begin{equation}
 z_i=\frac{q+T(w)}{12}-P_i,\qquad0\leq i<12.
 \label{iii:r25:inversa}
\end{equation}
Las doce relaciones \eqref{iii:r25:relaciones} y la relación
\eqref{iii:r25:cierre} son linealmente independientes sobre
$\mathbb Q$.
\end{theorem}
\begin{proof}
Para una salida de $\mathcal D$,
\[
 c_i-b_{i+1}
 =(z_i-z_{i+4})-(z_{i+1}-z_{i+4})=z_i-z_{i+1}.
\]
Tres incrementos consecutivos suman $b_i$ y la suma cíclica de
todos ellos es cero. Además, $P_i=z_0-z_i$, de donde
$q=12z_0-T(w)$. Se obtienen las condiciones y la fórmula inversa.
Recíprocamente, la congruencia hace enteras las coordenadas de
\eqref{iii:r25:inversa} y el cierre permite prolongarlas
cíclicamente. Sus diferencias adyacentes son $w_i$, luego $D_3z=b$
y
\[
 (D_4z)_i=w_i+w_{i+1}+w_{i+2}+w_{i+3}
 =w_i+b_{i+1}=c_i.
\]
La suma da $Q(z)=q$. La fórmula determina también la unicidad.
Finalmente, el cambio entero invertible
$(b,c)\mapsto(b,w=c-Sb)$ transforma las doce primeras relaciones
en $b=(I+S+S^2)w$: cada una despeja una coordenada diferente de
$b$. El cierre $\sum_iw_i=0$ añade una relación independiente.
\end{proof}

La imagen racional tiene dimensión $12$. La condición integral añade
la congruencia de orden $12$ de \eqref{iii:r25:congruencia}.
Los once incrementos $w_0,\ldots,w_{10}$ y la carga $q$ son
coordenadas suficientes, con $w_{11}=-\sum_{i=0}^{10}w_i$.
Las $25$ coordenadas conservan esos $12$ grados de libertad
mediante $13$ relaciones lineales. Si la terna proviene además de
\eqref{iii:r25:bloque}, la inversa pertenece a
$\{0,\ldots,999\}^{12}$ y recupera sus tres residuos por ventana.

\subsubsection{Concordancia con la transformación armónica}

La transformación $\mathcal H_{12}$ actúa mediante $H_4$ en las
órbitas $(r,r+3,r+6,r+9)$, $r=0,1,2$, donde
\[
 H_4=\begin{pmatrix}
 1&1&1&1\\1&1&-1&-1\\1&-1&1&-1\\1&-1&-1&1
 \end{pmatrix},\qquad H_4^2=4I_4.
\]
Para una terna compatible, sean
\[
 B_r=\sum_{j=0}^3c_{r+3j},\qquad
 a_0=\frac{q+2B_0+B_1}{3},\quad
 a_1=a_0-B_0,\quad a_2=a_1-B_1.
\]
El lector armónico tiene bloques
\[
 (\Pi_H(b,c,q))^{(r)}
 =(a_r,b_{r+3}-b_{r+9},b_r+b_{r+6},b_r-b_{r+6}).
\]

\begin{proposition}[Concordancia de reconstrucciones]
\label{iii:r25:hadamard}
Sobre $\operatorname{im}\mathcal D$ se cumple
\[
 \Pi_H\mathcal D=\mathcal H_{12},\qquad
 \mathcal D^{-1}=\tfrac14\mathcal H_{12}\Pi_H.
\]
La imagen de $\Pi_H$ restringida a ese dominio es
\[
 \mathcal L_H=
 \{u\in\mathbb Z^{12}:\mathcal H_{12}u\equiv0\pmod4\}.
\]
\end{proposition}
\begin{proof}
Para $z=\mathcal D^{-1}(b,c,q)$, sean
$h_r=\sum_{j=0}^3z_{r+3j}$. El desplazamiento de cuatro posiciones
lleva cada órbita a la siguiente; por tanto $B_r=h_r-h_{r+1}$.
La identidad $q=h_0+h_1+h_2$ da $a_r=h_r$.
En una órbita escribamos $(x_0,x_1,x_2,x_3)$ para sus coordenadas
y $d_j=x_j-x_{j+1}$. Se verifican
\[
 \begin{aligned}
 d_1-d_3&=x_0+x_1-x_2-x_3,\\
 d_0+d_2&=x_0-x_1+x_2-x_3,\\
 d_0-d_2&=x_0-x_1-x_2+x_3.
 \end{aligned}
\]
Son las tres filas restantes de $H_4x$. La inversión y la
caracterización integral se deducen de $\mathcal H_{12}^2=4I$.
\end{proof}

La condición de congruencia armónica controla la salida de $\Pi_H$;
la pertenencia de los canales a $\operatorname{im}\mathcal D$
requiere además \eqref{iii:r25:relaciones}--\eqref{iii:r25:congruencia}.
Un control negativo exacto consiste en tomar $b=0$, $q=0$ y
$c=e_0-e_3$. Las tres sumas orbitales de $c$ se anulan, de modo
que $\Pi_H(b,c,q)=0$, pero se incumplen las relaciones de imagen.
La lectura armónica no sustituye el control de compatibilidad de
los canales.

\subsection{Conjugación, covariancia y memoria del registro incidencial}
\label{iii:r25:transporte}

Sea $\rho(m)=13-m$. Si la conjugación de incidencias invierte el
orden de las ventanas, conserva las clases aritmética y de frontera
de las visitas e intercambia los canales de las trazas, entonces
\[
 (A_m^\dagger,C_m^\dagger,V_m^\dagger)
 =(A_{\rho(m)},-C_{\rho(m)},V_{\rho(m)}).
\]
Con $c_m=[C_m]_{10}$, la normalización decimal da
\begin{equation}
 z_m^\dagger=z_{\rho(m)}
 +100\mathbf1_{\{c_{\rho(m)}\ne0\}}-20c_{\rho(m)}.
 \label{iii:r25:conjugacion}
\end{equation}
En efecto, $[-C]_{10}=0$ si $[C]_{10}=0$ y vale
$10-[C]_{10}$ en los demás casos. La negación del canal se
aplica antes de reconstruir el bloque decimal.

La aplicación a una involución concreta conserva también los
extremos de las ventanas. Para una ruta
$\gamma=(x_0,\ldots,x_n)$ leída antes del avance,
\begin{equation}
 \sum_{k=0}^{n-1}f(x_{n-k})
 =\sum_{k=0}^{n-1}f(x_k)+f(x_n)-f(x_0).
 \label{iii:r25:extremos}
\end{equation}
Ambas sumas contienen los mismos términos interiores; la diferencia
es el extremo final menos el inicial. En rutas cerradas esa
corrección se anula. En ventanas abiertas se incorpora antes de
reducir módulo $10$.

\begin{proposition}[Consecuencia de una reducción periódica]
\label{iii:r25:periodo}
Si el observable de una familia fija de rutas satisface
$o(t+54)=o(t)$, sus multiplicidades se conservan bajo ese retorno
y el mismo lector local actúa en ventanas de nueve instantes,
entonces
\[
 (A,C,V)_{m+6}=(A,C,V)_m,\qquad z_{m+6}=z_m.
\]
\end{proposition}
\begin{proof}
La traslación $t\mapsto t+54$ lleva cada ventana biyectivamente
a la situada seis posiciones después y conserva cada incidencia
contada y su multiplicidad. La igualdad se transmite a los recuentos
y a sus normalizaciones decimales.
\end{proof}

Una presentación con doce bloques distintos conserva, por tanto,
información que no factoriza por ese observable de período $54$:
selección de incidencias, orientación, extremos, multiplicidad o
memoria. La proposición delimita el efecto de esa reducción; no
atribuye período $54$ a la historia completa.

\subsubsection{Determinación sobre órbitas y acumulación de eventos}

Fijar solamente $Q(z)=q$ deja una clase afín del retículo
\[
 \mathcal Z_0=\{v\in\mathbb Z^{12}:\sum_i v_i=0\}
 =\bigoplus_{i=0}^{10}\mathbb Z(e_i-e_{11}).
\]
Por ejemplo, $100\mathbf1+e_0$ y $100\mathbf1+e_1$ tienen
carga $1201$ y satisfacen los controles de imagen mediante sus
respectivos canales. Son testigos de las ecuaciones de compatibilidad,
sin atribución de pertenencia a las historias terminales HMT.

\begin{proposition}[Covariancia cíclica]
\label{iii:r25:covariancia}
Sea $\rho$ la traslación de ventanas de un registro $R$, de
período mínimo $p\mid12$, y sea $S$ la traslación de sus
coordenadas. Un lector covariante sobre esa órbita queda determinado
por $z\in\operatorname{Fix}(S^p)$. La condición $Q(z)=q$
admite soluciones enteras si y sólo si $12/p$ divide a $q$;
cuando existen, forman una clase afín de rango $p-1$.
\end{proposition}
\begin{proof}
La regla $F(\rho^jR)=S^jz$ está bien definida exactamente si
$S^pz=z$. Estos vectores repiten $12/p$ veces un bloque de
longitud $p$, y su carga es $\frac{12}{p}\sum_{i=0}^{p-1}z_i$.
La fórmula prueba la divisibilidad y el rango. La covariancia de
los canales se sigue de $D_aS=SD_a$ y $QS=Q$.
\end{proof}

En la carta armónica, la acción correspondiente es
$\widehat S_H=\mathcal H_{12}S\mathcal H_{12}/4$.
Si el origen está marcado, la acción transporta también esa marca.
El período de la órbita aquí indicada no se identifica con el de
la historia enriquecida. Para una lectura inicial $F_{108}$ ya
definida, la naturalidad bajo truncamiento se expresa mediante
$F_N=F_{108}\circ\tau_{N,108}$; esta identidad conserva la
evaluación inicial sin seleccionar su regla de producción.

\begin{proposition}[Composición aditiva de contribuciones incidenciales]
\label{iii:r25:eventos}
Supóngase que el estado inicial y cada evento $a_t$ producen,
mediante reglas previamente fijadas sobre el registro, contribuciones
$(\delta w_t,\delta q_t)$ tales que
\[
 \sum_i\delta w_{t,i}=0,
 \qquad \delta q_t+T(\delta w_t)\equiv0\pmod {12}.
\]
Con condiciones iniciales compatibles $(w^{(0)},q^{(0)})$,
la actualización por suma determina un único registro en cada etapa.
Su incremento es
\[
 \delta z_{t,i}=
 \frac{\delta q_t+T(\delta w_t)}{12}-P_i(\delta w_t).
\]
La construcción respeta la concatenación de segmentos y conserva
la lectura de un prefijo bajo prolongaciones posteriores.
\end{proposition}
\begin{proof}
Las condiciones de compatibilidad se conservan por suma y la
inversa \eqref{iii:r25:inversa} es aditiva en $(w,q)$.
La inducción fija cada estado. La suma sobre segmentos concatenados
es la suma de sus contribuciones; un prefijo recibe los mismos
sumandos antes y después de la prolongación.
\end{proof}

Esta proposición da un criterio suficiente para la acumulación
aditiva, no una exigencia de linealidad de todo lector TPK.
Cuando los eventos se transportan por acciones no triviales, sus
contribuciones se expresan en la carta terminal mediante el cociclo
afín correspondiente. Las reglas de evento y el estado inicial
conservan en ambos casos su procedencia incidencial.

El desarrollo anterior determina el paso
\[
 \mathscr L\longmapsto(A,C,V)\longmapsto z
 \longmapsto(D_3z,D_4z,Q(z))\longmapsto\mathcal H_{12}z
\]
con las condiciones de cada aplicación. Su dominio es distinto del
registro regional del apéndice~\ref{iii:nucleo-finito}: allí las
ecuaciones $X_aL_a=Y_a$ determinan los levantamientos a partir de
transiciones ternarias. Aquí las incidencias de visitas y trazas
determinan bloques enteros y sus canales. La unicidad de una
reconstrucción no reemplaza la producción del registro que utiliza
la otra. Ambas conservan sus primitivas, operaciones y pruebas en
la cadena APP--TRIT--TPK.
